\section{A product-torus tail homotopy in the actual two-colour complex}
\label{sec:r12-product-torus-tail-homotopy}

\paragraph{Status and scope.}
This source is submitted as \textbf{PROVED THIS ROUND}, pending
independent review of its exact bytes.  The construction contributors
are \texttt{/root} and \texttt{/root/r11\_challenge}.  Its new interface
is a restricted, explicitly represented integration operator on actual
forms, followed by a finite actual homotopy.  It is not an arbitrary
proper analytic pushforward theorem, not a descent theorem, and not
a homotopy on every analytic family on a torus.  The analytic base
below is an ordinary open in a smooth space over the original field
$K$; no calculus on a nonreduced norm-parameter base is used.

\subsection{Objects, equivalence and the precise parameter restriction}

Let $Y$ be a smooth pure-dimensional good $K$-analytic space without
boundary.  Put $\mathcal T=\mathbf G_{m,K}^{\mathrm{an}}$, let $t$
be its coordinate, let $\tau=v(t)$, and write
$\pi:Y\times\mathcal T\to Y$ and $i_1:Y\to Y\times\mathcal T$
for projection and the section $t=1$.  All absolute values and
weights have their original-$K$ normalization.

A form $u\in B(Y\times\mathcal T)$ is \emph{product-presented on
bounded parameter slabs} if for each $y\in Y$ and each compact
interval $I\subset\mathbb R$ there are an ordinary neighborhood
$V$ of $y$, base analytic units $f=(f_1,\ldots,f_a)$ on $V$, and
a finite native $\delta$-form $\gamma$ such that
\[
 u|_{V\times\{\tau\in I\}}=t_{(f,t)}^\star\gamma.
 \tag{PTH1}
\]
The equality is required on a neighborhood of the displayed slab.
One may restrict the native form to an ambient open neighborhood
of its image and use a finite partition there.  Enlarging the unit
tuple and subdividing is permitted.  Units $t^b f_j$, $b\in\mathbb Z$,
are permitted because their valuations are integral affine functions
of $(v(f),\tau)$.  General units depending on $t$, such as $1-t$,
are \emph{not} permitted unless a separate presentation (PTH1) is
actually supplied on the whole slab in question.

Equality and zero mean equality and zero in Mihatsch's actual
$B$ on the whole analytic product, not merely equality on one
chosen tropicalization.  Thus this is a subsheaf of actual forms,
not a free native coefficient space or its completion.  Finite sums,
wedge products, and both full differentials preserve this class.
An analytic change of base preserves it by actual pullback.

We call such a form \emph{locally slab-supported} if every $y$ has
an ordinary neighborhood $V$ and a compact interval $I$ such that
\[
                 \operatorname{Supp}(u|_{V\times\mathcal T})
                        \subseteq\{\tau\in I\}.
 \tag{PTH2}
\]
We call it \emph{positive-tail zero} if every $y$ has such a $V$
and a constant $M$ for which $u=0$ on $V\times\{\tau>M\}$.
These are local conditions on the \emph{whole actual form}; they
do not assert that a chosen native expression vanishes coefficient
by coefficient outside a slab.

\subsection{Finite native integration and affine base change}

Let $p_a:\mathbb R^a\times\mathbb R\to\mathbb R^a$ be projection,
with standard multiplicity one in the last direction.  For a native
$\delta$-form $\gamma$ supported in a compact $\tau$-slab, define
$p_{a,*}\gamma$ by finite polyhedral fiber integration.  In a
subordinate finite complex, if
$\gamma=\sum_\sigma\alpha_\sigma\wedge[\sigma,\mu_\sigma]$,
its formula is
\[
 p_{a,*}\gamma=
 \sum_\sigma
 p_{a,\mu_\sigma/\nu_\sigma,*}\alpha_\sigma
                  \wedge[p_a(\sigma),\nu_\sigma].
 \tag{PTH3}
\]
Here $\nu_\sigma$ is any positive weight on the image, and the
quotient weight is the one determined by the exact sequence of
linear spans.  Each coefficient is an actual finite fiber integral,
not an abstract dual functional.  Polyhedral subdivision gives
piecewise smooth coefficients.  Affine pushforward commutes with
$d'$ and $d''$ as currents, so both derivatives of (PTH3) are
polyhedral.  Therefore (PTH3) is a native $\delta$-form.  Its total
bidegree is lower by $(1,1)$, irrespective of the dimensions of
the individual carriers.  This uses the definition and affine
functoriality of native $\delta$-forms, not a pushforward on an
arbitrary tropical or analytic space.

\begin{lemma}[Affine product base change]
For every affine map $a:\mathbb R^b\to\mathbb R^a$ and every
such $\gamma$,
\[
 a^*p_{a,*}\gamma
       =p_{b,*}(a\times\mathrm{id})^*\gamma.
 \tag{PTH4}
\]
All pullbacks in this identity are the native graph pullbacks.
\end{lemma}

\begin{proof}
Factor $a$ as a surjective affine map followed by an affine closed
embedding.  In the surjective case, linear changes of coordinates
make it a coordinate projection.  Its pullback is product with the
standard fiber space, and (PTH4) follows directly from (PTH3) by
Fubini, with the product quotient weights.

For an embedding $j$, put $J=j\times\mathrm{id}$ and let
$C=j_*1$ be its weighted affine image cycle.  The native graph
pullback and projection formula give
\[
 \begin{split}
 j_*p_{b,*}J^*\gamma
 &=p_{a,*}J_*J^*\gamma
   =p_{a,*}(p_a^*C\wedge\gamma)\\
 &=C\wedge p_{a,*}\gamma
   =j_*j^*p_{a,*}\gamma.
 \end{split}
 \tag{PTH5}
\]
Pushforward by an affine embedding is injective on polyhedral
currents: a smooth compactly supported test form on the image
extends to an ambient tubular neighborhood.  Cancelling this
injective map proves (PTH4).  Functoriality gives the general case.
The invoked native projection and graph formulas are Theorem 4.1
and Proposition/Definition 4.2 of
\texttt{arXiv:2107.12067v3}.  No non-surjective pullback for an
arbitrary tropical space has been used.
\end{proof}

\subsection{The product-cycle test and actual zero relations}

We use the following original-field product fact.  If $K_0\subset Y$
is an affinoid domain, $g$ is any finite tuple of its analytic units,
and $A_I$ is a closed standard annulus with valuation interval $I$,
then, away from the analytic boundaries,
\[
 T(K_0\times A_I,(g,t))=T(K_0,g)\times[I,1].
 \tag{PTH6}
\]
Here $K_0$ is an affinoid domain, not an extension of the ground
field.  A direct verification also fixes the normalization.  Above
a point of $\Sigma(K_0,g)$ the maximal-dimensional skeleton for
$(g,t)$ consists of the Gauss extensions
\[
 \left|\sum_j a_jt^j\right|_{x,\tau}
                 =\max_j |a_j(x)|\exp(-j\tau).
 \tag{PTH7}
\]
Indeed maximal graded transcendence degree forces the reduction
of $t$ to be transcendental over the graded residue field of the
base point; this is precisely the Gauss extension.  Conversely
(PTH7) has the required extra transcendence degree.  Thus the
weighted skeleton is the product skeleton.  Locally, a finite flat
toric coordinate projection of base degree $e$ becomes, after
adjoining $t$, a coordinate projection of the same degree $e$.
The parameter coordinate is primitive and contributes weight one.
This proves (PTH6), also with nonrational radii.  It is the
one-center product case of the skeleton computation in the proof
of Proposition 3.21 of \texttt{arXiv:2112.10018v4}.  The boundary
images to exclude are the base boundary image times $I$ and the
two annulus endpoints; we will choose $I$ larger than the parameter
support.  Dimension-deficient base images contribute zero on both
sides and cause no change to this argument.

\begin{lemma}[Entire actual zero preservation]
If $t_{(f,t)}^\star\gamma=0$ on $V\times\mathcal T$ and $\gamma$
is supported in a compact parameter slab, then
\[
                       t_f^\star(p_{a,*}\gamma)=0
                                   \quad\hbox{in }B(V).
 \tag{PTH8}
\]
\end{lemma}

\begin{proof}
Let $K_0\subset V$ be \emph{any} affinoid domain and let
$g$ be \emph{any} unit refinement of $f|_{K_0}$, say
$f|_{K_0}=a\circ g$ on valuation coordinates.  Choose a standard
annulus whose interval strictly contains the parameter support.
The forward implication of the actual zero criterion, Lemma 4.4
of \texttt{arXiv:2112.10018v4}, applied to this affinoid product
and the refinement $(g,t)$, gives
\[
 (T(K_0,g)\times\mathbb R)
                \wedge(a\times\mathrm{id})^*\gamma=0
 \tag{PTH9}
\]
away from the base boundary image.  The annulus endpoints make
no contribution because of the strictly smaller support slab.
Apply the native parameter projection.  The native projection
formula and (PTH4) turn (PTH9) into
\[
                T(K_0,g)\wedge a^*p_{a,*}\gamma=0
 \tag{PTH10}
\]
away from that boundary image.  Since $K_0$ and $g$ were arbitrary,
the reverse implication of the same Lemma 4.4 proves (PTH8).

This checks the entire actual relation kernel, not just the
initial chart, Gauss chambers, separated exterior products, or a
probe topology.  Additional units depending on $t$ are not being
declared unnecessary for testing input zero: input zero is
\emph{already an actual equality}.  Product refinements suffice
to deduce the required equalities for \emph{all base refinements}
of the proposed output.
\end{proof}

\begin{theorem}[Restricted actual product integration]
There is a well-defined linear operator on locally slab-supported,
product-presented actual forms,
\[
 \Pi:B^{p,q}_{\mathrm{prod,slab}}(Y\times\mathcal T)
                                  \longrightarrow B^{p-1,q-1}(Y),
 \qquad
 \Pi(t_{(f,t)}^\star\gamma)=t_f^\star p_{a,*}\gamma.
 \tag{PTH11}
\]
It commutes with restriction to ordinary base opens, with both
full differentials, and with analytic base change whenever both
spaces satisfy the stated hypotheses.  It satisfies the actual
projection formula with every actual base form.  No claim is
made for forms not supplied with the product presentation.
\end{theorem}

\begin{proof}
First work on a base neighborhood with one support slab.  Choose
a smooth scalar $\chi(\tau)$, equal to one near that slab and
zero outside a slightly larger compact slab.  Choose (PTH1) on
a neighborhood of the larger slab, multiply the native expression
by $\chi$, and extend it by zero beyond the slab.  In the native
base variables one may work locally and then use a finite smooth
partition subordinate to the ambient presentation neighborhoods.
Formula (PTH3) supplies a finite native output.

For two presentations use a common base unit refinement and
pull back their native expressions to it.  Formula (PTH4) compares
the two native outputs on that common chart.  The difference of
the two cut-off inputs is actual zero, so (PTH8) makes the
outputs equal in actual $B$.  The same argument proves
independence of the cutoff and of ambient subdivisions and
partitions.  On ordinary base overlaps it proves equality of
the outputs, which therefore glue as actual forms.  Notice that
this construction does not infer native support from actual
support: the cutoff makes native support proper explicitly.

Native pushforward commutes with $d'$ and $d''$.  Cutoff derivative
terms are actual zero because their support is outside the
actual support of the form; (PTH8) kills their outputs.  This
proves differential compatibility.  For a base form, include
its finitely many presenting units in the common chart and
use the native projection formula.  Finally analytic base
change replaces the tuple $f$ by its actual analytic pullback
and leaves the displayed native integral unchanged.  The actual
pullback theorem is Theorem 4.5 of the same v4 source.  Thus
all the asserted compatibilities hold as actual identities.
\end{proof}

For clarity about topology, fix a finite unit tuple, finite
carrier complexes, a compact parameter slab and compact base
coefficient sets.  On the corresponding finite smooth coefficient
Fr\'echet spaces, (PTH3) is continuous into one fixed finite
image subdivision.  This follows by differentiating the finitely
many fiber integrals with affine endpoints; each target seminorm
is bounded by finitely many source seminorms of bounded order.
The quotient by entire actual zero relations is permitted by
(PTH8).  One may therefore give the represented source the
coefficient-final topology of these specified stages and obtain
a continuous map into the actual coefficient-final target.
This statement neither identifies that source topology with
the topology inherited from all $B(Y\times\mathcal T)$ nor
asserts a closed range, strictness, Hausdorff stage comparison,
or a full OPEN-KERNEL result.  No completion occurs in (PTH11).

\subsection{The section identity and the homotopy}

Put $h=\tau_+=\max(\tau,0)$.  Its native double derivative is
the multiplicity-one hyperplane $\{\tau=0\}$.

\begin{lemma}[Section identity in exactly the product class]
For every product-presented actual $u$,
\[
                             \Pi(Dh\wedge u)=i_1^*u.
 \tag{PTH12}
\]
\end{lemma}

\begin{proof}
The left side is locally slab-supported, and only a neighborhood
of the zero slab is needed.  Present $u=t_{(f,t)}^\star\gamma$
there.  Let $j_0:\mathbb R^a\to\mathbb R^a\times\mathbb R$
be the affine section at zero.  In native $\delta$-forms,
the full graph-pullback projection formula gives
\[
 p_{a,*}(D\tau_+\wedge\gamma)
       =p_{a,*}(j_{0,*}1\wedge\gamma)
       =p_{a,*}j_{0,*}j_0^*\gamma=j_0^*\gamma.
 \tag{PTH13}
\]
This is valid for all carriers, including normal first-form
components and corner terms; it is not a restriction of their
coefficients by hand.  On the other hand, the genuine analytic
section $t=1$ takes the entire unit presentation $(f,t)$ to
$(f,1)$, so actual pullback gives
$i_1^*u=t_f^\star j_0^*\gamma$.  This proves (PTH12).
For a general analytic family the last sentence has no such
native product description, and (PTH12) need not hold.
\end{proof}

Use the exact array complex and the operators $J_k$ from HCT1--6:
\[
 L_k^s=\bigoplus_{i=0}^k B^{s-i,s-k+i},\qquad
 (\partial_k u)_i=d''u_i+d'u_{i+1}.
 \tag{PTH14}
\]
For a positive-tail-zero product-presented $u\in L_k^s$, define
\[
                         \mathcal H_k u=\Pi J_k(h,u).
 \tag{PTH15}
\]
Every component of $J_k(h,u)$ vanishes on $\tau<0$, and on
the given positive tail because $u$ and its derivatives vanish
there.  Thus (PTH15) is a legitimate locally slab-supported
actual integration, not a divergent formal integral.  Its
degree is
\[
                        \mathcal H_k u\in L_{k+1}^s(Y).
 \tag{PTH16}
\]
In particular for a closed triple in the original order
$q=(A,B,M)$, (PTH15) uses the literal QHT expression
\[
             K_h(q)=(d''h\wedge A,-d'h\wedge B,
                         -h\wedge d''M,-d'h\wedge M),
 \tag{PTH17}
\]
with its original quadruple ordering and no discarded mixed term.

\begin{theorem}[Actual tail homotopy, with its nonclosed defect]
For every $k\ge1$ and every such $u$,
\[
 \partial_{k+1}\mathcal H_k u
        =i_1^*u-\mathcal H_{k-1}(\partial_k u).
 \tag{PTH18}
\]
Consequently an independently verified closed family gives the
explicit actual primitive
\[
               i_1^*u=\partial_{k+1}\mathcal H_k u.
 \tag{PTH19}
\]
At degree zero the precise endpoint is
\[
             \partial_1\mathcal H_0 g
                         =i_1^*g-\Pi(h\wedge Dg).
 \tag{PTH20}
\]
\end{theorem}

\begin{proof}
Apply (PTH11) to the entire HCT5 identity
$\partial J_k(h,u)=Dh\wedge u-J_{k-1}(h,\partial u)$
and then use (PTH12).  Both terms are locally slab-supported;
$\partial u$ inherits the same zero tail.  Equivalently one may
insert a native cutoff equal to one on $[0,M]$ and with its
derivatives in a strictly larger positive tail.  All derivative
terms of that cutoff are actual zero, not merely negligibly
small, and vanish under (PTH11).  HCT6 proves the endpoint.
No assertion that closedness by itself implies a primitive has
entered the argument: the genuine product family and its entire
uniform zero tail are explicit additional data.
\end{proof}

\subsection{Pressure tests and the external-row application interface}

The base unit tuple is unrestricted: it may include
$(T_1,T_2,T_1+T_2)$ on its unit locus, or the units defining a
relative curve with moving punctures and any number of passive
directions.  Their actual zero relations are retained because
(PTH8) tests every base affinoid refinement.  Nothing translates
these base coordinates independently or replaces their whole
relation kernel by product probes.  Real intercept walls and
mixed parameter jets are part of the finite native integral.
Nonrational analytic points participate in (PTH7) and in the
actual zero criterion.  No field extension, rational-point
density argument, division by a domain degree, or unproved
high-order actual descent is used.

For the closed external-row family $\mathcal T^t$ of EFT2--4,
once that family's exact finite product-presentation and tail
certificate is independently verified, the precise substitution
in (PTH19) is
\[
 H(S,v;g)=\Pi K_{\tau_+}(\mathcal T^t),\qquad
              \partial_3H(S,v;g)=\mathcal T(S,v;g).
 \tag{PTH21}
\]
This would close the specific ECR10 primitive requirement with
zero extending error, for this specified defect only.  It does
not demand a strict witness for ASP's diagonal symbol cycles.
The construction of the whole family is essential; an arbitrary
closed actual triple has not been given such a family here.
Further Gaussian chart triangles still need their literal
pair/triple comparisons and the frozen RAG/RSP hypotheses before
one can claim an actual Green form for a new original cycle.

\paragraph{Exact inputs and primary-source audit.}
HCT:
\path{27af6c36a7819874f2a3137bf4ee6db4b8f8d4f7ffd99fc37d419b2144790549}.
QHT:
\path{aa10132dc6e8217dfa4c896699c76eaed90ce4d632abc5e0147cb9e278c423e0}.
The native formulas used above were read in
\texttt{arXiv:2107.12067v3} (28 August 2022), Sections 2.3 and 3,
Theorem 4.1 and Proposition/Definition 4.2.  The analytic interfaces
were read in \texttt{arXiv:2112.10018v4} (7 September 2024),
Proposition 2.17, the product case in Proposition 3.21,
Definition 4.2, Lemma 4.4 and Theorem 4.5.  Example 2.45(3)
of the latter explicitly warns against an unrestricted
$\delta$-form pushforward on tropical spaces.  The finite ambient
native construction and all-refinement argument above are the
additional restricted bridge, not an assertion that this warning
has been removed in general.
